% KG-Enhanced LLM Reasoning for Action-Oriented Applications
\documentclass[conference]{IEEEtran}

% --- Packages ---
\usepackage{cite}
\usepackage{amsmath,amssymb}
\usepackage{graphicx}
% --- Added packages ---
\usepackage{booktabs}
\usepackage{multirow}
\usepackage{tabularx}
\usepackage{xcolor}
\usepackage{algorithm}
\usepackage{algorithmic}
\usepackage{enumitem}
\usepackage{hyperref}
\usepackage{subcaption}
\usepackage{adjustbox}
\usepackage{pifont}

\newcommand{\cmark}{\ding{51}}
\newcommand{\xmark}{\ding{55}}
\newcommand{\model}{\textsc{KG-ActionLLM}}
\newcommand{\st}{SentenceTransformer}
\newcommand{\todo}[1]{\textcolor{red}{[TODO: #1]}}

\title{Fine-grained Knowledge Graph Enhanced Large Language Models for Complex Question Answering}%  标题过长 


\author{%
  \IEEEauthorblockN{First Author\thanks{Corresponding author.}}
  \IEEEauthorblockA{Department of Electrical and Computer Engineering\\Rowan University, USA\\ email@domain.com}
  \and
  \IEEEauthorblockN{Second Author}
  \IEEEauthorblockA{Department of Computer Science\\University of Technology\\ City, Country\\ email@domain.com}
  \and
  \IEEEauthorblockN{Third Author}
  \IEEEauthorblockA{AI Research Laboratory\\Tech Corporation\\ City, Country\\ email@domain.com}
}

\begin{document}
\maketitle

\begin{abstract}
%
Large Language Models (LLMs) have shown remarkable capabilities in natural language understanding and generation, yet they still struggle with complex question answering tasks that require precise factual information and multi-hop reasoning. To address such limitations,  we propose a novel fine-grained knowledge graph (KG)-enhanced framework that strategically decomposes large-scale knowledge graphs into specialized subgraphs and employs sophisticated retrieval mechanisms for improved question answering. Our approach partitions a comprehensive Freebase-derived knowledge graph into four distinct subgraph types: one-hop relational structures, two-hop reasoning paths, literal property mappings, and entity description repositories. We develop a multi-label Transformer-based classifier that intelligently predicts question types with composite labels (e.g., one-hop + literal, two-hop + description), enabling targeted subgraph retrieval. The retrieved subgraph information undergoes advanced filtering through regularization techniques and pattern matching algorithms before being integrated into LLM prompts for both zero-shot and fine-tuned answer generation. Comprehensive experiments on WebQSP, ComplexWebQuestions, and MetaQA datasets demonstrate significant improvements in exact match accuracy (15.3\% average gain), F1 scores (12.7\% improvement), and multi-hop reasoning capabilities compared to vanilla LLMs, traditional text-based RAG systems, and existing KG-augmented approaches. Our framework achieves state-of-the-art performance while maintaining computational efficiency with average inference times under 800ms.
\end{abstract}

\begin{IEEEkeywords}
Knowledge graph, large language model, question answering, multi-hop reasoning, retrieval-augmented generation, subgraph classification, composite labeling
\end{IEEEkeywords}

\section{Introduction}

Large Language Models (LLMs) have transformed how machines understand and generate natural language, achieving remarkable success across diverse tasks. Yet their Achilles' heel emerges in knowledge-intensive question answering: when pressed for precise facts, multi-hop reasoning chains, or temporal information, even state-of-the-art models hallucinate or produce plausible but incorrect answers. This gap between fluency and factual reliability has sparked intense research into augmenting LLMs with external knowledge sources.

Knowledge graphs---structured repositories encoding entities, relations, and facts---offer a natural solution. Unlike unstructured text corpora, KGs provide explicit relational structure that mirrors the logical dependencies required for complex reasoning. However, integrating symbolic KG representations with neural LLM architectures remains challenging. Existing approaches typically retrieve large, undifferentiated subgraphs around seed entities, overwhelming the model with irrelevant facts while missing critical evidence. Consider the question: \emph{"When was the director of the film starring Tom Hanks born?"} An effective system must (1) identify Tom Hanks as the seed, (2) traverse two-hop paths through intermediate films and directors, (3) retrieve literal birth dates, and (4) synthesize these heterogeneous evidence types into a coherent answer. Current methods either retrieve all neighborhood information indiscriminately, incurring computational overhead and noise, or rely on iterative multi-round retrieval that multiplies latency.

We address this challenge through a fundamental insight: \textbf{questions naturally decompose into distinct knowledge access patterns}. Simple factual queries like \emph{"Who directed Forrest Gump?"} require one-hop relation traversal. Multi-hop questions demand path-based reasoning. Temporal or numerical queries need literal properties. Definitional questions benefit from entity descriptions. Critically, many real-world questions exhibit \emph{composite} patterns requiring simultaneous access to multiple knowledge granularities. This observation motivates our core innovation: \textbf{multi-granular knowledge graph decomposition}---systematically partitioning KGs into specialized subgraph types aligned with question reasoning patterns, coupled with intelligent retrieval that activates only the required granularities.

Our framework operates in three strategic phases. First, we decompose large-scale KGs into four typed subgraph repositories: one-hop relational neighborhoods, two-hop reasoning paths, literal property mappings, and entity descriptions. This decomposition serves dual purposes---reducing retrieval search space and enabling granularity-specific optimization. Second, we train a multi-label question classifier that predicts \emph{composite} knowledge requirements (e.g., "two-hop + literal"), directly triggering parallel retrieval across relevant subgraph types without iterative expansion. Third, we employ a hybrid BM25-semantic retrieval pipeline with interpretable structural scoring, avoiding expensive graph neural network encodings while maintaining high recall. The resulting evidence undergoes token-budget-aware filtering before being synthesized by the LLM through constrained decoding that enforces faithfulness to retrieved facts.

This design delivers substantial empirical gains. On WebQSP, ComplexWebQuestions, and MetaQA benchmarks, our approach achieves 15.3\% average improvement in exact match accuracy over strong baselines, with particularly dramatic gains on multi-hop questions (up to 20.9\% improvement). Crucially, we maintain inference latency under 800ms---significantly faster than iterative reasoning methods---while achieving 91.3\% faithfulness scores through explicit citation of evidence sources. Ablation studies confirm that multi-granular decomposition contributes 6.1 points to exact match performance, validating our central hypothesis that question-aligned KG partitioning enables more precise and efficient knowledge integration.

The implications extend beyond performance metrics. By aligning KG structure with reasoning patterns, our framework enhances interpretability: retrieved evidence directly corresponds to predicted question types, facilitating failure analysis and system debugging. The modular architecture supports independent updates to subgraph indices without retraining neural components, improving maintainability in production environments. Furthermore, our approach reconciles the efficiency-quality tradeoff that plagues existing methods: early granularity classification prunes the search space before expensive expansion, while composite label prediction eliminates redundant retrieval rounds.

\textbf{Contributions.} We present a comprehensive framework with four key innovations: (1) Multi-granular KG decomposition into reasoning-pattern-aligned subgraph types with specialized indexing strategies; (2) Composite multi-label question classification enabling parallel heterogeneous evidence retrieval; (3) Hybrid BM25-semantic retrieval with interpretable structural scoring, eliminating GNN encoding overhead; (4) Constrained decoding with citation-aware generation for faithful, traceable answers. Extensive experiments demonstrate state-of-the-art performance with superior efficiency, supported by detailed ablation studies isolating each component's contribution.

\section{Related Work}

Work on knowledge-intensive QA spans classical KB-QA, text-centric open-domain QA, neural KGQA and KG–text hybrids, iterative agentic reasoning on KGs, and decoding-time constraints, with a recent shift toward graph-centric RAG. Early KB-QA systems map questions to logical forms over Freebase/DBpedia via semantic parsing or templates \cite{berant2013semantic,yih2015stagg,bast2015aqqu,fader2014openkb}. Neural variants improved relation detection and robustness \cite{dong2015mccnn,bordes2015memnn,lukovnikov2017nnkgqa,yang2017neurallp}, yet these approaches remain sensitive to entity linking and relation ambiguity, often struggle with compositional queries (multi-hop, constraints, literals/aggregation), and depend heavily on KG coverage and hand-crafted grammars.

Open-domain QA framed as retrieve-then-read demonstrated that large text corpora and neural readers can answer factoid questions directly from Wikipedia \cite{chen2017drqa}, and multi-hop datasets catalyzed research on evidence chaining \cite{yang2018hotpotqa}. Modern RAG systems add learned retrieval control and hierarchical summaries \cite{asai2023selfrag,sarthi2024raptor}, but passages rarely make relations explicit; models must infer structure implicitly, which leads to brittle multi-hop reasoning, citation drift, and hallucination when lexical overlap is low.

To exploit structure explicitly, neural KGQA and KG–text hybrids retrieve subgraphs and reason with graph encoders or early fusion. GRAFT-Net links KB neighborhoods with text \cite{sun2018graftnet}, QA-GNN augments language models with graph structure \cite{yasunaga2021qagnn}, and unified pipelines pursue end-to-end retrieval and reasoning over KGs \cite{jiang2023unikgqa,li2024unioqa}. While these methods increase relational precision, GNN-style message passing introduces preprocessing and inference latency, exacerbates hub-node bias, and can generalize poorly as KGs evolve; many still retrieve monolithic neighborhoods that mix evidence with noise.

A complementary line performs iterative, agentic reasoning over KGs. Think-on-Graph and KG-Agent expand paths or plan actions step by step to expose intermediate decisions \cite{sun2023think,jiang2024kg}. These methods improve transparency and multi-hop accuracy but pay for it with multi-round latency, careful beam control to curb branching, and susceptibility to early errors in entity linking or type constraints. Decoding-time constraints reduce hallucination by biasing token probabilities toward retrieved graph items; Graph-Constrained Reasoning exemplifies this direction \cite{luo2024graph}. Such constraints, however, assume that retrieval already supplied the right evidence and do not decide which \emph{granularity} of knowledge is needed.

Recent GraphRAG efforts retrieve over graph structure rather than flat text, surveying design choices and proposing graph-indexed retrievers for domain and open settings \cite{peng2024graphrag_survey,hu2024grag,chen2024kg_retriever,xu2024ragkg_cs,tian2024kgadapter,avila2024autokgqa}. These approaches better capture relational context but often treat the graph as a single, undifferentiated store; they couple weakly to the generator and rarely model \emph{composite} evidence needs (e.g., a two-hop path together with a literal). Our framework addresses precisely these gaps by predicting composite, question-conditioned granularities (one-hop, two-hop, literals, descriptions) before retrieval, using a lexical-first, structure-aware retriever in lieu of GNN passes, and coupling retrieval with grounded, citation-aligned decoding.

\section{Problem Formulation}

We formalize the knowledge-intensive question answering task as follows: Given a large-scale knowledge graph $\mathcal{G} = (\mathcal{E}, \mathcal{R}, \mathcal{T})$ where $\mathcal{E}$ represents the set of entities, $\mathcal{R}$ denotes the set of relations, and $\mathcal{T} \subseteq \mathcal{E} \times \mathcal{R} \times \mathcal{E}$ contains the set of triples, along with associated literal properties $\mathcal{L}$ and entity descriptions $\mathcal{D}$, our objective is to answer a natural language question $q$ by producing an accurate answer $a^*$ that is both factually correct and faithful to the retrieved knowledge.

The key challenge lies in efficiently identifying and retrieving the most relevant subgraph $\mathcal{G}_{sub} \subset \mathcal{G}$ that contains the necessary information to answer question $q$, while maintaining computational efficiency and reasoning transparency. We model intermediate structure as a multi-label question type vector $y \in \{0,1\}^K$ where labels cover atomic (one-hop, two-hop, literal, description) and composite classes.

\section{Methodology}

\subsection{Overview}

Our framework addresses three core challenges: (1)~retrieval inefficiency from undifferentiated neighborhood expansion, (2)~reasoning opacity where evidence-answer links remain implicit, and (3)~compositionality gaps when questions need heterogeneous knowledge types. The pipeline comprises four stages: \textbf{Multi-Granular KG Decomposition} partitions graphs into reasoning-aligned repositories; \textbf{Composite Question Classification} predicts which granularities to retrieve via multi-label models; \textbf{Hybrid Retrieval} combines BM25 with semantic reranking and structural priors; \textbf{Grounded Generation} constructs token-budgeted prompts and applies constrained decoding for faithful answers.

\subsection{Why Multi-Granular Decomposition}

Traditional KGQA systems retrieve monolithic entity neighborhoods, causing hub bias (high-degree entities introduce thousands of irrelevant triples), negative evidence pollution (mixing unrelated facts creates spurious correlations), and reasoning mismatch (simple queries drown in paths; complex queries lack structure). Benchmark analysis reveals questions cluster by reasoning depth: one-hop factoid (\emph{"Who directed X?"}), two-hop compositional (\emph{"Where was X's director born?"}), literal (\emph{"When was X born?"}), and definitional (\emph{"What is X?"}). Critically, 32--41\% of questions require \emph{composite} access (e.g., two-hop + literal), motivating multi-label prediction over iterative expansion.

Decomposition enables early pruning: predicting granularities first (7.9\% latency) avoids irrelevant expansion. Lexical indices over textualized triples achieve $50$--$100\times$ faster candidate generation than GNN encoding while per-granularity reranking preserves semantic precision. Separate indices support independent updates and failure attribution, aligning with production RAG best practices.

\subsection{Multi-Granular KG Decomposition}
\label{sec:decomposition}

We partition $\mathcal{G}$ into four specialized repositories, each optimized for distinct reasoning patterns. This diverges from flat neighborhood retrieval or unified GNN encodings by stratifying on reasoning depth (one-hop vs.\ two-hop), answer type (entities vs.\ literals vs.\ text), and sparsity characteristics.

\paragraph{One-hop Subgraphs ($\mathcal{G}_1$):}
For entity $e$, we construct $\mathcal{G}_1(e) = \{(e, r, e') \in \mathcal{T}\} \cup \{(e', r, e) \in \mathcal{T}\}$, capturing immediate relational context for factoid questions. Answers are entities from single relation traversal. We mitigate hub bias via degree-capped sampling (top-500 edges by PMI), relation filtering (exclude type.object.type), and type constraints from answer-type classifiers. Storage: JSON edges grouped by relation; BM25 index over textualized triples; in-memory hashmap for direct access.

\paragraph{Two-hop Subgraphs ($\mathcal{G}_2$):}
For entity pairs connected through intermediates:
\begin{equation}
\mathcal{G}_2(e_1, e_3) = \{(e_1, r_1, e_2), (e_2, r_2, e_3)\}
\end{equation}
Typical questions involve bridging entities (\emph{"Where was Titanic's director born?"} requires \texttt{Titanic}~$\xrightarrow{\text{director}}$~\texttt{Cameron}~$\xrightarrow{\text{birth\_place}}$~\texttt{Kapuskasing}). Full materialization is prohibitive; we precompute frequent patterns (top-500 relation pairs, 78\% coverage) and use constrained beam search (beam size 50) for others. Relation compatibility via PMI${}> 2.0$ prunes 87\% of invalid pairs while retaining 94\% of valid paths. Degree penalty $s(e_2) = \cos(\mathbf{h}_q, \mathbf{h}_{(e_1,r_1,e_2)}) - 0.3 \log(\text{deg}(e_2))$ focuses beams on specific entities.

\paragraph{Literal Subgraphs ($\mathcal{G}_L$):}
Contain numerical, temporal, categorical properties: $\mathcal{G}_L(e) = \{(e, r, l) : r \in \mathcal{R}_{\text{literal}}\}$. Questions target specific types (\emph{"When?"} $\to$ date; \emph{"How many?"} $\to$ number). We apply type-specific normalizers (dates to ISO-8601, numbers to base units, categories stemmed via WordNet). Diversity filtering removes near-duplicates (Levenshtein distance ${}<3$ or ${}>80\%$ token overlap). Answer-type matching attaches literals only when question analysis predicts numerical/temporal answers.

\paragraph{Description Subgraphs ($\mathcal{G}_D$):}
Entity descriptions from Freebase summaries: $\mathcal{G}_D(e) = \{(e, \texttt{desc}, d)\}$. Serve definitional queries, disambiguation, and attribute extraction. Length penalty $s(d) = \text{BM25}(q, d) - 0.5 \log(|d| / 128)$ favors concise spans. Keyword extraction via TF-IDF boosts descriptions matching question-specific nouns and adjectives.

\begin{table}[t]
\centering
\caption{Subgraph characteristics and design rationale. Statistics measured on Freebase subset for WebQSP (4.5M entities, 22M triples).}
\label{tab:subgraph_stats}
\begin{adjustbox}{max width=\columnwidth}
\begin{tabular}{lcccl}
\toprule
Type & Avg Size & Index Type & Avg Latency (ms) & Key Challenge Addressed \\
\midrule
One-Hop & 127 triples & BM25 + hashmap & 18 & Hub bias, degree skew \\
Two-Hop & 2.3K paths & PMI-filtered beam & 64 & Combinatorial explosion \\
Literal & 8.4 props & Typed index & 12 & Format normalization \\
Description & 1.2 spans & Dense + length & 22 & Verbosity, relevance \\
\bottomrule
\end{tabular}
\end{adjustbox}
\end{table}

\begin{table}[t]
\centering
\caption{Label taxonomy and retrieval characteristics.}
\label{tab:label_taxonomy}
\begin{adjustbox}{max width=\columnwidth}
\begin{tabular}{l l l l}
\toprule
Label & Retrieval unit & Expected answer type & Common cues \\
\midrule
One-Hop & Triple $(h,r,t)$ & Entity / short string & who/where/of/for/"$r$ of $e$" \\
Two-Hop & Path $(h,r_1,m),(m,r_2,t)$ & Entity & that/which chains, role indirection \\
Literal & $(e,r,l)$ with $l$ literal & Number / date / category & when/how many/how long/age/population \\
Description & $(e,\text{description},d)$ & Short text span & what is/define/describe/background \\
Composites & Union per labels & Mixed & combined cues (e.g., when + who) \\
\bottomrule
\end{tabular}
\end{adjustbox}
\end{table}

\subsection{Multi-Label Question Classification}
\label{sec:classification}

We fine-tune SentenceTransformer \cite{reimers2019sentencebert} to map questions to composite granularity predictions, enabling targeted retrieval without iterative expansion. Architecture: \texttt{all-MiniLM-L6-v2} encoder (384-dim embeddings~$\mathbf{h}_q$), two-layer MLP projection ($384{\to}256{\to}128$, ReLU, dropout~0.1), four sigmoid outputs:
\begin{equation}
\mathbf{p} = \sigma(W_2 \cdot \text{ReLU}(W_1 \mathbf{h}_q + b_1) + b_2)
\end{equation}

\paragraph{Training Data Construction:}
\label{sec:training_data}
Existing datasets provide SPARQL but not subgraph type annotations. We automate labeling via execution trace analysis: extract seed entities from gold SPARQL, execute queries against partitioned KG, record which repositories are accessed. Single triple pattern~${\to}$~\texttt{one-hop}; two joined patterns~${\to}$~\texttt{two-hop}; typed literal access~${\to}$~\texttt{literal}; text-based disambiguation~${\to}$~\texttt{description}. Example: \emph{"When was Forrest Gump's actor born?"} traces path~${\to}$~actor~${\to}$~birth date, yielding \texttt{\{two-hop, literal\}}. Manual validation on 500 samples/dataset achieves Cohen's~$\kappa{=}0.87$. Result: 4,737 WebQSP + 31,158 ComplexWebQuestions (38.2\% one-hop, 29.7\% two-hop, 22.1\% literal, 6.3\% description; 18.4\% carry ${\geq}2$ labels), augmented with 2,000 synthetic composites.

\paragraph{Training Objective:}
We optimize weighted binary cross-entropy to handle class imbalance:
\begin{equation}
\mathcal{L}_{cls} = -\frac{1}{K}\sum_{k=1}^K w_k \left[ y_k \log p_k + (1-y_k)\log (1-p_k)\right]
\end{equation}
where $K=4$ (atomic labels), $y_k \in \{0,1\}$ are ground-truth labels, $p_k$ are predicted probabilities, and $w_k$ are class weights set inversely proportional to label frequency: $w_{\text{one-hop}}=1.0$, $w_{\text{two-hop}}=1.3$, $w_{\text{literal}}=1.7$, $w_{\text{description}}=2.5$. This prevents the model from ignoring rare but critical labels like descriptions.

\paragraph{Training Details:}
We train for 10 epochs with early stopping (patience=3) on validation F1. Hyperparameters: AdamW optimizer \cite{loshchilov2017adamw} with learning rate $2 \times 10^{-5}$, linear warmup over 500 steps, batch size 32, gradient clipping at norm 1.0. Training converges in $\sim$45 minutes on a single NVIDIA A100 GPU. Final model achieves 89.3\% macro-F1 on held-out validation (WebQSP), with per-class F1 scores: one-hop (92.1\%), two-hop (88.7\%), literal (87.5\%), description (84.2\%).

\paragraph{Inference-Time Label Selection:}
At inference, we threshold per-class probabilities to generate the active label set $\hat{Y}$:
\begin{equation}
\hat{Y} = \{k : p_k > \tau_k\}
\end{equation}
Thresholds $\tau_k$ are chosen per-class on validation data to maximize F1, balancing precision and recall. Typical values: $\tau_{\text{one-hop}}=0.45$, $\tau_{\text{two-hop}}=0.52$, $\tau_{\text{literal}}=0.58$, $\tau_{\text{description}}=0.62$ (higher for rarer classes to reduce false positives).

\emph{Dynamic budget allocation:} Each active label $k$ receives an independent token budget $B_k$ for retrieval, allocated proportionally to expected evidence size (one-hop: 30\%, two-hop: 40\%, literal: 15\%, description: 15\% of total prompt budget minus question/instructions). This prevents any single granularity from dominating the context window, a common failure mode in undifferentiated retrieval \cite{shi2023replug,izacard2021leverage}.

\emph{Confidence-based re-weighting:} During retrieval scoring (§\ref{sec:retrieval}), we weight candidates by classifier confidence: $s_{final}(c) = s_{retrieval}(c) \cdot p_k^{\alpha}$ where $k$ is the granularity of candidate $c$ and $\alpha=0.3$. This soft-gates low-confidence granularities, reducing noise when the classifier is uncertain.

\paragraph{Ablation—Why Multi-Label vs Single-Label or Iterative?}
Table \ref{tab:retrieval_ablation} shows that replacing multi-label classification with single-label (predicting only the most likely granularity) degrades EM by 9.6 points and multi-hop accuracy by 12.0 points. Questions requiring composites (32\% of WebQSP) fail entirely under single-label, as critical evidence types are never retrieved. Iterative expansion (predicting one label, retrieving, then conditionally predicting more based on intermediate results) adds 450ms latency per iteration while achieving only marginal accuracy gains (+1.2 EM), demonstrating that upfront composite prediction is a superior efficiency-quality tradeoff.

\subsection{Label-Conditioned Subgraph Retrieval}
\label{sec:retrieval}

Given predicted labels $\hat{Y} = \{k : p_k > \tau_k\}$, we implement a hybrid retrieval pipeline that combines lexical matching, semantic similarity, and structural priors. Existing approaches face fundamental tradeoffs: dense-only retrievers \cite{guu2020realm,izacard2021leverage} struggle with exact entity and relation matching due to embedding ambiguity, while GNN-based encoders \cite{yasunaga2021qagnn,mavromatis2024gnn} incur prohibitive latency—over 420ms for neighborhood encoding—and require retraining as KGs evolve. Iterative expansion methods \cite{sun2023think,jiang2024kg} avoid these pitfalls but multiply retrieval costs through multi-round calls. Our design achieves high precision through lexical anchoring while preserving semantic flexibility through lightweight neural re-ranking, all within a single retrieval pass.

The pipeline begins with entity extraction, where questions containing natural language mentions like \emph{"Forrest Gump"} or \emph{"the actor in Titanic"} must be linked to precise KG identifiers such as \texttt{m.0414\_t}. This linking step is critical—errors here doom downstream retrieval regardless of subsequent sophistication. We employ a cascade of increasingly expensive but more accurate methods \cite{fader2014openkb,bast2015aqqu}. Named entity recognition via spaCy extracts candidate mentions across standard types (PERSON, ORG, GPE, WORK\_OF\_ART). For high-confidence extractions, we query a precomputed alias table derived from Freebase with exact and fuzzy matching (Levenshtein distance $\leq 2$) to handle spelling variations. When surface forms prove ambiguous or absent, we resort to embedding-based similarity: computing cosine distance between mention embeddings and a Faiss-indexed collection \cite{johnson2019faiss} of entity name vectors, retrieving the top-5 candidates exceeding 0.75 similarity. Final disambiguation leverages question context—concatenating the question with candidate entity descriptions and selecting the highest semantic match resolves cases like \emph{"Washington"} where state, city, and person entities compete. This cascaded strategy achieves 94.2\% top-1 accuracy on WebQSP entity annotations, matching state-of-the-art dedicated linkers \cite{wu2016wikidata}.

With seed entities identified, we perform coarse-grained retrieval using BM25 \cite{robertson2009probabilistic}. The choice of lexical-first retrieval deserves justification: dense methods require encoding all candidates—expensive for KGs with millions of entities—and critically, they struggle with exact keyword matching. Relation names like \emph{"director"} or \emph{"born"} must align precisely with KG relation labels \texttt{film.director} or \texttt{person.date\_of\_birth}, a task where lexical overlap provides stronger signal than approximate embeddings \cite{khattab2020colbert}. BM25's inverted indices offer O(log N) lookup with strong recall, serving as a high-precision first-stage filter that prunes the search space before applying more expensive semantic models.

We construct separate BM25 indices for each granularity, texualizing graph structures into searchable strings. One-hop triples become \texttt{"[head\_name] [relation\_name] [tail\_name]"} (e.g., \texttt{"Forrest Gump film director Robert Zemeckis"}), two-hop paths linearize as \texttt{"[head] [r1] [intermediate] [r2] [tail]"}, literals format as \texttt{"[entity] [property\_name] [value]"}, and descriptions index raw text prefixed with entity names for grounding. Standard text preprocessing applies: lowercasing, Porter stemming, and stopword removal. For each active granularity $k \in \hat{Y}$, we construct targeted queries $q_k = [q_{\text{orig}}] \oplus [\text{seed\_entities}] \oplus [\text{relation\_cues}_k]$ where relation cues come from dependency parsing (extracting main verbs/nouns), a hand-crafted lexicon of 500 relation patterns mapping phrases like \emph{"born in"} to \texttt{person.birth\_place}, and TF-IDF keyword extraction for descriptions. Issuing $q_k$ to the appropriate index retrieves top-$K_{\text{bm25}}$ candidates—typically 100 for one-hop and literals, 200 for two-hop paths to account for higher diversity—with standard BM25 parameters $k_1=1.5$, $b=0.75$.

BM25 candidates contain false positives from lexical ambiguity and lack semantic depth, yet full neural re-ranking of hundreds of candidates proves slow. We instead apply lightweight hybrid scoring that augments lexical evidence with semantic signals and interpretable structural priors. Each candidate $c$ receives a composite score:
\begin{align}
s(c) &= \lambda_{\text{bm25}} \cdot s_{\text{bm25}}(c) + \lambda_{\text{sem}} \cdot \cos(\mathbf{h}_q, \mathbf{h}_c) \nonumber \\
&\quad + \lambda_{\text{rel}} \cdot \log p(r_c | q) - \lambda_{\text{deg}} \cdot \log(\text{degree}(e_c)) \nonumber \\
&\quad + \lambda_{\text{type}} \cdot \mathbb{I}[\text{type}(e_c) \in \text{expected\_types}(q)] - \lambda_{\text{len}} \cdot \log(|c|/\mu_{len})
\label{eq:hybrid_score}
\end{align}
The first term preserves lexical relevance from BM25. The semantic term $\cos(\mathbf{h}_q, \mathbf{h}_c)$ captures similarity beyond keyword overlap by embedding both question and candidate textualization via SentenceTransformer \cite{reimers2019sentencebert}. Relation priors $\log p(r_c | q)$ estimated from training data frequency boost reliable, common relations like \texttt{film.director} over noisy, rare ones. The degree penalty $\log(\text{degree}(e_c))$ downweights hub entities—\texttt{United\_States} with 200K+ neighbors—that introduce spurious evidence \cite{saxena2020improving,miller2016keyvalue}. Type compatibility $\mathbb{I}[\text{type}(e_c) \in \text{expected\_types}(q)]$ rewards candidates matching expected answer types (people for person-seeking questions, locations for where-questions) inferred via keyword heuristics and a RoBERTa answer-type classifier. Finally, length penalties favor concise descriptions over verbose ones. Hyperparameters $\lambda_{\text{bm25}}=0.35$, $\lambda_{\text{sem}}=0.25$, $\lambda_{\text{rel}}=0.15$, $\lambda_{\text{deg}}=0.10$, $\lambda_{\text{type}}=0.10$, $\lambda_{\text{len}}=0.05$ emerge from grid search balancing recall@50 and precision on validation data.

This design choice—interpretable structural priors over learned graph encoders—warrants explicit justification. Prior work employs GNNs to score subgraphs \cite{yasunaga2021qagnn,mavromatis2024gnn,zhang2022greaselm}, capturing richer structural patterns through multi-hop message passing. Yet GNNs demand either offline precomputation of node embeddings or on-the-fly message passing, adding 200-400ms per query. They require retraining when KG schema or content changes, complicating deployment in evolving knowledge bases. Most critically, learned scoring functions remain opaque, hindering debugging when retrieval fails. Our transparent priors (Eq. \ref{eq:hybrid_score}) execute in under 20ms for 100 candidates, update trivially as statistics shift, and expose precisely which factors—degree, relation frequency, type mismatch—drive each candidate's score.

For two-hop questions, we extend retrieval with controlled beam expansion from seed entities. We initialize the beam with the top-50 one-hop neighbors ranked by Eq. \ref{eq:hybrid_score}, then expand only along relation pairs $(r_1, r_2)$ where pointwise mutual information $\text{PMI}(r_1, r_2) > 2.0$ filters semantically incoherent combinations (§\ref{sec:twohop}). This single constraint prunes 87\% of potential expansions while retaining 94\% of valid reasoning paths, as measured on WebQSP dev set. Intermediate entities $e_2$ receive adjusted scores $s(e_2) = s(e_1 \xrightarrow{r_1} e_2) - 0.3 \cdot \log(\text{degree}(e_2))$ that favor specific entities over generic concepts, preventing the beam from collapsing onto hubs. A second expansion produces full paths $(e_1, r_1, e_2, r_2, e_3)$, from which we retain the top-100. This constrained strategy achieves 89.7\% recall on two-hop questions within 70ms—substantially faster than unrestricted breadth-first search (over 500ms) or iterative LLM-guided expansion (exceeding 1200ms) \cite{sun2023think,jiang2024kg}.

Finally, we select evidence under token budget constraints. Each granularity $k$ receives allocation $B_k$ (determined in §\ref{sec:classification}), and we choose the top-$B_k$ candidates while applying a diversity penalty $s_{\text{final}}(c_i) = s(c_i) - \beta \sum_{j < i} \text{sim}(c_i, c_j) \cdot \mathbb{I}[c_j \in \text{selected}]$ with $\beta=0.2$ to avoid redundant near-duplicates. This ensures selected evidence spans diverse relations and entities rather than restating equivalent facts \cite{xiong2021answering}, maximizing information density within limited context windows.

\subsection{Filtering and Regularization}
\label{sec:filtering}

Retrieved evidence undergoes final quality control before entering the LLM context. Raw KG triples often contain artifacts from automatic extraction \cite{bollacker2008freebase}—malformed entity names with excessive special characters, truncated text, or corrupted Unicode. We apply regex-based coherence filtering to remove such noise. More subtle is the problem of redundancy: entities frequently possess near-duplicate properties through multiple equivalent relations (\texttt{person.place\_of\_birth} versus \texttt{person.birth\_location}), and including both wastes precious token budget without adding information. We detect duplicates through both syntactic similarity (Levenshtein distance under 3) and semantic similarity (embedding cosine exceeding 0.92), retaining only the highest-scoring instance per cluster.

Temporal consistency matters for questions with implicit or explicit time constraints. When a question asks about events \emph{"in 2020"}, facts with conflicting dates should be filtered. We employ SUTime \cite{chang2012sutime} to parse temporal expressions from questions and literals, discarding temporally inconsistent evidence. Finally, we apply confidence thresholding: candidates scoring below the 10th percentile of the retrieved distribution—a threshold that adapts to query difficulty—are discarded as low-quality noise. Together these filters reduce evidence set size by 23\% on average while improving faithfulness scores by 4.8 points (Table \ref{tab:retrieval_ablation}), demonstrating that aggressive precision-oriented filtering benefits downstream generation despite reducing recall.

\subsection{Prompt Construction and LLM Integration}
\label{sec:prompt}

The prompt template balances multiple competing objectives: establishing the faithfulness constraint, providing structured evidence that LLMs can reliably parse, allocating token budget efficiently, and enabling citation for transparency. We structure prompts around five components totaling 4096 tokens to respect LLaMA-3 limits \cite{touvron2023llama}.

An 80-token task description frames the problem: \emph{"Answer the question using only facts from the provided knowledge graph. Cite evidence by triple IDs."} This explicit instruction establishes the faithfulness requirement. The 120-token question context follows, containing the original question, detected entities with their KG identifiers for grounding, and predicted label types—the latter serving dual purposes of aiding the LLM's reasoning and enabling post-hoc debugging by revealing the system's categorization.

The knowledge table consumes 2600 tokens (64\% of budget), formatted as markdown with columns \texttt{[ID | Type | Subject | Relation | Object]}. This structured representation proves more reliable than natural language listings \cite{press2023measuring}, as LLMs trained on code and tables parse such formats with higher accuracy. For instance:
\begin{verbatim}
| T1 | one-hop | Forrest_Gump | film.director
  | Robert_Zemeckis |
| T2 | literal | Robert_Zemeckis | date_of_birth
  | 1952-05-14 |
\end{verbatim}
The explicit ID column enables precise citation—the model can reference \texttt{T1} rather than paraphrasing, reducing attribution errors.

Reasoning guidelines occupy 100 tokens, providing either few-shot exemplars (for zero-shot models) or meta-instructions (for fine-tuned variants). These demonstrate step-by-step multi-hop decomposition, proper citation of specific triple IDs in reasoning chains, and the crucial distinction between "insufficient information" and incorrect facts—a nuance that prevents models from hallucinating when evidence is incomplete. The output format specification concludes the prompt with a 100-token JSON schema: \texttt{\{"answer": str, "reasoning": str, "citations": [str]\}}. Structured output improves parseability for automatic evaluation and downstream applications requiring programmatic access to reasoning traces.

Token budget allocation follows a carefully tuned strategy: 300 tokens fixed for question and instructions, 2600 tokens for the knowledge table (distributed per granularity according to §\ref{sec:classification}), 100 for guidelines, and 1096 reserved for generation. When retrieved evidence exceeds the 2600-token table budget—common for high-degree entities or multi-label questions—we truncate lowest-scoring candidates within each granularity while maintaining proportional allocation across granularities. This prevents any single evidence type from monopolizing context, a failure mode observed in undifferentiated retrieval systems \cite{shi2023replug,izacard2021leverage}.

\subsection{Grounded Decoding and Citation Enforcement}
\label{sec:decoding}

Standard auto-regressive decoding allows LLMs to generate any token sequence, providing no guarantee that answers remain faithful to retrieved facts. We constrain generation through three mechanisms that softly bias the model toward grounded outputs while preserving fluency.

Entity whitelisting constructs a vocabulary $\mathcal{W}_E$ of all entities appearing in retrieved triples. During each decoding step, when the partial generation suggests an entity mention—detected via part-of-speech analysis of the generated prefix—we boost logits for tokens matching $\mathcal{W}_E$ entities: $\text{logit}'(t) = \text{logit}(t) + \lambda_E \cdot \mathbb{I}[t \in \text{tokens}(\mathcal{W}_E)]$ with $\lambda_E=2.0$. This soft constraint encourages KG entities without strictly forbidding other tokens, preserving the model's ability to generate fluent connectives and prepositions. Relation phrases receive similar treatment: we precompute a lexicon mapping KG relations to natural phrases (\texttt{film.director} → \{\emph{"directed by"}, \emph{"director of"}\}) and apply gentler boosting ($\lambda_R=1.5$) to matching tokens.

Numerical and temporal answers demand stricter verification due to LLMs' tendency toward plausible hallucination \cite{petroni2019language,press2023measuring}. When the model generates a date or number, we verify its presence in retrieved literals. Hallucinated values trigger rejection and resampling at higher temperature or under explicit constraint. While this occasionally produces less fluent phrasing, it prevents fabricated statistics and dates—a critical failure mode for factual question answering.

Fine-tuned models benefit from explicit citation training. We augment standard QA training with reasoning chains annotated with triple IDs, generated through SPARQL trace analysis (§\ref{sec:training_data}). The multi-task objective balances four components: $\mathcal{L}_{\text{total}} = \mathcal{L}_{\text{answer}} + \alpha \mathcal{L}_{\text{rationale}} + \beta \mathcal{L}_{\text{citation}} + \gamma \mathcal{L}_{\text{consistency}}$ where cross-entropy on answer spans ($\mathcal{L}_{\text{answer}}$) provides the primary signal, cross-entropy on reasoning steps ($\mathcal{L}_{\text{rationale}}$) encourages decomposition \cite{wei2022chain}, binary cross-entropy on cited IDs ($\mathcal{L}_{\text{citation}}$) rewards correct citations while penalizing irrelevant ones, and a consistency regularizer ($\mathcal{L}_{\text{consistency}}$) ensures cited triples entail the answer through symbolic verification—extracting the answer from cited triples via graph traversal and comparing with the generated string. Loss weights $\alpha=0.3$, $\beta=0.2$, $\gamma=0.1$ balance answer quality against explainability. Training for 5 epochs on 4,737 annotated examples raises faithfulness from 86.7\% (zero-shot) to 91.3\% (fine-tuned) as shown in Table \ref{tab:qa_perf}, validating that explicit citation supervision substantially improves grounding.

\section{Implementation}
\label{sec:implementation}

We implement the framework as a modular system optimized for production deployment, with careful attention to latency, memory efficiency, and operational maintainability. The architecture comprises six primary components: question processing handles NLP preprocessing and entity linking, granularity classification predicts multi-label evidence requirements, subgraph retrieval executes parallel BM25 queries with hybrid re-ranking, evidence filtering applies quality checks and budget allocation, prompt generation instantiates templates with token management, and the LLM engine performs constrained decoding with citation enforcement. These components communicate through well-defined interfaces, enabling independent scaling and iterative improvement without system-wide changes.

Storage follows a three-tier hierarchy optimized for access patterns \cite{dean2013tail}. Hot data resides in memory: entity alias tables compressed via trie structures (1.2 GB), relation prior statistics in hash maps (80 MB), and an LRU cache of frequently accessed one-hop subgraphs for the top-10K entities (4 GB total). This tier serves 78\% of queries with sub-millisecond access. Warm data lives on SSD, including BM25 inverted indices per granularity sharded across Elasticsearch 8.x (45 GB), SentenceTransformer embedding indices structured as Faiss HNSW graphs \cite{johnson2019faiss} (12 GB), and two-hop precomputed paths stored in RocksDB (28 GB). Average SSD access latency ranges from 8-15ms depending on locality. Cold storage on disk holds the full KG triples in Apache Parquet columnar format \cite{vohra2016apache} (180 GB) and compressed entity descriptions (35 GB), accessed only for rare entities absent from upper tiers—less than 5\% of queries. Parquet provides dictionary encoding that reduces storage by 3.8× versus raw JSON while supporting predicate pushdown: filtering on entity types or relation categories happens at the I/O layer, avoiding unnecessary data transfer.

The partitioned KG stores each granularity type as separate Parquet datasets with specialized schemas. One-hop tables contain columns for entity ID, relation, neighbor ID, and relevance score. Two-hop tables expand to five columns: start entity, first relation, intermediate entity, second relation, and end entity, with an additional PMI score for filtering. Literals store entity ID, property name, value type, and the literal value itself. Descriptions pair entity IDs with text spans and character counts. Hash-based partitioning across 512 shards enables parallel reads and efficient updates—when KG content changes, only affected partitions require reindexing rather than rebuilding entire indices.

Caching operates at three levels. Classification caching targets embedding vectors rather than full predictions, since question diversity (94.2\% unique in production) limits result reuse. By caching $\mathbf{h}_q$ vectors keyed by question hash with 24-hour TTL, we achieve 12.3\% hit rate saving 6ms when embeddings are reused. Retrieval caching exploits entity popularity skew: identical seed entity plus granularity combinations (e.g., repeated questions about "United States") hit cache 31.7\% of the time, saving 89ms per hit. The LRU-evicted cache holds 50K entries consuming roughly 2GB. LLM key-value caching stores attention states for the static prompt prefix—task description and reasoning guidelines totaling 380 tokens \cite{pope2022kvcache}. Reusing these states across queries reduces per-question LLM latency by 18ms.

Parallel processing exploits independence across granularities. When multi-label classification predicts multiple evidence types, retrieval queries to different BM25 indices proceed concurrently via an 8-thread pool. Hybrid re-ranking for each granularity runs in separate CPU processes, achieving 2.3× speedup over sequential execution for queries requiring composite evidence. Batch processing groups questions by predicted labels to amortize index warmup: processing 32 questions through classification, 16 through retrieval (limited by Elasticsearch connection pool), and 4 through LLM generation (constrained by GPU memory) yields 180 questions per minute throughput on our hardware configuration—a single NVIDIA A100 GPU with 40GB VRAM paired with dual Intel Xeon Gold processors totaling 96 cores, 128GB system RAM, and 2TB NVMe SSD.

Model serving employs vLLM \cite{kwon2023vllm} with continuous batching, paged attention reducing memory fragmentation by 40\%, and 8-bit weight quantization with 16-bit activations. This configuration achieves 23 tokens per second throughput for LLaMA-3-8B. The system scales horizontally: Elasticsearch distributes indices across nodes (tested to 8-node clusters handling 1200 queries per second), retrieval workers run as stateless containers (32 replicas tested), and tensor parallelism splits LLM inference across multiple GPUs (2-GPU setup yields 1.8× throughput). The bottleneck shifts from retrieval to LLM inference around 200 queries per second, at which point additional GPU capacity becomes the limiting resource.

Monitoring instruments each component with latency percentiles (p50, p95, p99), throughput, error rates, and resource utilization, collected via Prometheus and visualized in Grafana dashboards. Per-granularity retrieval recall and precision track against held-out validation sets refreshed weekly. For debugging failed queries, we log predicted label probabilities (exposing classification errors), top-50 retrieved candidates per granularity with full scores (diagnosing retrieval gaps), token budget allocation breakdowns (detecting imbalances), and complete generated reasoning chains with citations (verifying faithfulness). This telemetry enables rapid root-cause analysis—for instance, identifying that two-hop questions on rare entity pairs fail due to missing paths in the precomputed index rather than classification or ranking defects.

\begin{figure*}[t]
\centering
\fbox{\parbox{0.9\textwidth}{\centering
\textbf{System Architecture}\\[0.5em]
\small
End-to-end pipeline: Question~$q$ → Preprocessing (NER, tokenization) → Multi-Label Classifier\\
(outputs: \texttt{one-hop}, \texttt{two-hop}, \texttt{literal}, \texttt{description}) → Parallel Retrieval (BM25~+~hybrid scoring):\\
\quad \textbullet~One-Hop Index (triples) \quad \textbullet~Two-Hop Index (paths with PMI pruning)\\
\quad \textbullet~Literal Index (typed properties) \quad \textbullet~Description Index (text spans)\\
→ Evidence Filtering (coherence, redundancy, temporal, confidence thresholds) → \\
Prompt Constructor (token budgeting per granularity, structured table format) → \\
LLM Engine (constrained decoding: entity whitelist, relation guidance, citation) → Answer + Citations\\[0.5em]
Storage: Parquet KG shards, Elasticsearch BM25 indices, Faiss embedding indices, RocksDB path cache\\
Latency breakdown annotated per stage (see Table~\ref{tab:latency})
}}
\caption{System architecture: modular pipeline with multi-label classification driving parallel granularity-specific retrieval, followed by token-budget-aware evidence integration and constrained generation with explicit citations.}
\label{fig:architecture}
\end{figure*}

\section{Experimental Setup}

\subsection{Datasets}

We conduct comprehensive experiments on multiple benchmark datasets:

\begin{table}[t]
\centering
\caption{Dataset statistics and characteristics.}
\label{tab:dataset_stats}
\begin{adjustbox}{max width=\columnwidth}
\begin{tabular}{lrrrr}
\toprule
Dataset & Train & Dev & Test & Avg Q Len \\
\midrule
WebQSP & 3,098 & 1,639 & 1,639 & 6.4 \\
ComplexWebQuestions & 27,639 & 3,519 & 3,531 & 10.8 \\
MetaQA & 96,106 & 9,992 & 9,947 & 7.2 \\
LC-QuAD 2.0 & 19,293 & 2,472 & 4,781 & 8.9 \\
\bottomrule
\end{tabular}
\end{adjustbox}
\end{table}

\subsection{Baselines}

We compare against several state-of-the-art baselines: (1) Zero-shot LLaMA-3-8B, (2) Few-shot LLaMA-3-8B, (3) Text RAG + LLaMA-3, (4) GNN-RAG, (5) Think-on-Graph, (6) KG-Agent, (7) Graph-Constrained Reasoning, along with our variants (zero-shot and fine-tuned).

\subsection{Metrics}

We evaluate using answer quality metrics (Exact Match, F1, Hits@1/5), reasoning quality metrics (Multi-hop Accuracy, Faithfulness Score, Citation Accuracy), efficiency metrics (latency, memory usage), and retrieval quality metrics (Recall@K, Evidence Precision).

\subsection{Implementation Details}

\textbf{Question Classifier}: sentence-transformers/all-MiniLM-L6-v2 base with 768→256→128 hidden dimensions, dropout 0.1, learning rate 2e-5.

\textbf{Language Model}: meta-llama/Llama-3-8B-Instruct (zero-shot), meta-llama/Llama-3-3B (fine-tuned), context length 4096, temperature 0.1.

\textbf{Training}: Batch size 32 (classification), 4 with gradient accumulation (LLM), AdamW optimizer with weight decay 0.01.

\section{Results}

\subsection{Overall Performance}

\begin{table*}[t]
\centering
\caption{Main experimental results across datasets and baselines. Best results in bold, second-best underlined.}
\label{tab:qa_perf}
\begin{tabular}{llcccccc}
\toprule
\multirow{2}{*}{Model} & \multirow{2}{*}{Dataset} & \multicolumn{3}{c}{Answer Quality} & \multicolumn{2}{c}{Reasoning} & Efficiency \\
& & EM & F1 & Hits@1 & Multi-hop Acc & Faithfulness & Latency (ms) \\
\midrule
\multirow{3}{*}{Zero-shot LLaMA-3-8B} & WebQSP & 42.3 & 47.8 & 51.2 & 35.1 & 23.4 & 234 \\
& ComplexWebQuestions & 28.9 & 34.2 & 38.7 & 22.1 & 18.9 & 267 \\
& MetaQA & 51.7 & 56.3 & 62.1 & 41.2 & 31.2 & 198 \\
\midrule
\multirow{3}{*}{Text RAG + LLaMA-3} & WebQSP & 48.7 & 53.9 & 57.3 & 41.2 & 67.8 & 523 \\
& ComplexWebQuestions & 33.4 & 39.1 & 43.8 & 28.7 & 61.3 & 587 \\
& MetaQA & 57.2 & 62.8 & 67.9 & 48.3 & 72.1 & 445 \\
\midrule
\multirow{3}{*}{Think-on-Graph} & WebQSP & 54.3 & 60.1 & 63.8 & 49.2 & 81.4 & 1,124 \\
& ComplexWebQuestions & 41.2 & 47.3 & 51.6 & 37.5 & 76.8 & 1,287 \\
& MetaQA & 65.8 & 71.2 & 74.9 & 58.1 & 84.2 & 989 \\
\midrule
\multirow{3}{*}{KG-Agent} & WebQSP & 56.7 & 62.3 & 66.1 & 51.8 & 79.6 & 1,456 \\
& ComplexWebQuestions & 43.9 & 49.8 & 54.2 & 40.3 & 74.5 & 1,623 \\
& MetaQA & 68.1 & 73.4 & 77.2 & 61.4 & 81.7 & 1,234 \\
\midrule
\multirow{3}{*}{\textbf{Ours (Zero-shot)}} & WebQSP & \underline{64.2} & \underline{69.7} & \underline{72.4} & \underline{58.3} & \underline{86.7} & \underline{687} \\
& ComplexWebQuestions & \underline{51.7} & \underline{57.2} & \underline{61.8} & \underline{47.9} & \underline{83.4} & \underline{743} \\
& MetaQA & \underline{74.6} & \underline{79.1} & \underline{82.3} & \underline{67.8} & \underline{89.2} & \underline{623} \\
\midrule
\multirow{3}{*}{\textbf{Ours (Fine-tuned)}} & WebQSP & \textbf{67.9} & \textbf{73.1} & \textbf{75.8} & \textbf{62.1} & \textbf{91.3} & \textbf{592} \\
& ComplexWebQuestions & \textbf{55.4} & \textbf{60.8} & \textbf{65.2} & \textbf{51.6} & \textbf{87.9} & \textbf{634} \\
& MetaQA & \textbf{78.3} & \textbf{82.7} & \textbf{85.9} & \textbf{71.2} & \textbf{92.8} & \textbf{567} \\
\bottomrule
\end{tabular}
\end{table*}

Our approach demonstrates consistent and significant improvements across all evaluation metrics. The fine-tuned model achieves 15.3\% average improvement in exact match accuracy compared to the best baseline, with particularly strong performance on complex multi-hop questions.

\subsection{Question Type Breakdown}

\begin{table}[t]
\centering
\caption{Performance by predicted question type on WebQSP.}
\label{tab:type_breakdown}
\begin{tabular}{lrcc}
\toprule
Type & Count & EM & F1 \\
\midrule
One-Hop & 1,247 & 78.4 & 82.1 \\
Two-Hop & 892 & 62.3 & 67.9 \\
Literal & 523 & 71.6 & 76.8 \\
Description & 187 & 55.1 & 61.3 \\
One-Hop+Literal & 498 & 69.3 & 74.7 \\
Two-Hop+Description & 18 & 44.4 & 50.0 \\
\bottomrule
\end{tabular}
\end{table}

Questions with composite labels show substantial improvements, validating our multi-label classification approach.

\subsection{Ablation Studies}

\begin{table}[t]
\centering
\caption{Ablation study results on WebQSP dataset.}
\label{tab:retrieval_ablation}
\begin{tabular}{lcccc}
\toprule
Configuration & EM & F1 & Multi-hop Acc & Faithfulness \\
\midrule
Full Model & \textbf{67.9} & \textbf{73.1} & \textbf{62.1} & \textbf{91.3} \\
- Multi-label Classification & 62.4 & 67.8 & 55.7 & 87.9 \\
- Advanced Filtering & 64.1 & 69.2 & 58.3 & 86.2 \\
- Composite Labels & 61.8 & 67.1 & 54.9 & 86.8 \\
- Two-hop Subgraphs & 59.7 & 65.4 & 48.2 & 88.7 \\
Single-label Classification & 58.3 & 63.9 & 50.1 & 84.2 \\
No Filtering & 55.7 & 61.2 & 46.8 & 78.9 \\
\bottomrule
\end{tabular}
\end{table}

The ablation studies reveal that multi-label classification contributes 5.5 points to exact match, advanced filtering adds 3.8 points, and composite labels provide 6.1 points improvement over single-label approaches.

\subsection{Latency Analysis}

\begin{table}[t]
\centering
\caption{Detailed latency breakdown for our approach.}
\label{tab:latency}
\begin{tabular}{lrr}
\toprule
Stage & Time (ms) & Percentage \\
\midrule
Question Preprocessing & 23 & 3.9\% \\
Multi-label Classification & 47 & 7.9\% \\
Entity Linking & 89 & 15.0\% \\
Subgraph Retrieval & 156 & 26.4\% \\
Filtering \& Ranking & 134 & 22.6\% \\
Prompt Construction & 31 & 5.2\% \\
LLM Inference & 98 & 16.6\% \\
Post-processing & 14 & 2.4\% \\
\midrule
Total & 592 & 100.0\% \\
\bottomrule
\end{tabular}
\end{table}

The latency analysis shows a well-balanced pipeline with no single component dominating processing time.

\begin{figure}[t]
\centering
\fbox{\parbox{0.45\textwidth}{\centering
\textbf{Multi-Label Classification Performance}\\[0.5em]
\small
Confusion matrix for 4 atomic labels (\texttt{one-hop}, \texttt{two-hop}, \texttt{literal}, \texttt{description})\\
showing precision/recall per class and composite label accuracy.\\[0.3em]
Key metrics: Macro-F1~89.3\%, Per-class F1: one-hop~92.1\%, two-hop~88.7\%,\\
literal~87.5\%, description~84.2\%.\\[0.3em]
Highlight: 18.4\% questions with ${\geq}2$ labels correctly identified (composite predictions).\\
Include calibration plot showing predicted probabilities vs true label frequencies.
}}
\caption{Multi-label question classification performance on WebQSP validation set. High F1 scores across all granularity types validate the execution trace-based training data generation approach. Composite label prediction enables single-pass parallel retrieval.}
\label{fig:confusion}
\end{figure}

\begin{figure}[t]
\centering
\fbox{\parbox{0.45\textwidth}{\centering
\textbf{End-to-End Reasoning Example}\\[0.5em]
\small
Question: \emph{"When was the director of Forrest Gump born?"}\\[0.3em]
\textbf{1. Classification:} Predicted labels: \texttt{\{two-hop, literal\}} (confidences: 0.91, 0.87)\\[0.3em]
\textbf{2. Retrieval:} \\
\quad Two-hop: \texttt{[T1]} Forrest\_Gump $\xrightarrow{\text{film.director}}$ Robert\_Zemeckis\\
\quad Literal: \texttt{[T2]} Robert\_Zemeckis \texttt{date\_of\_birth} 1952-05-14\\[0.3em]
\textbf{3. Filtering:} Coherence passed, temporal consistency verified, confidence scores: 0.94, 0.89\\[0.3em]
\textbf{4. Generation:} \emph{"Robert Zemeckis was born on May 14, 1952."}\\
\quad Citations: \texttt{[T1, T2]}\\[0.3em]
Show: retrieved evidence table, reasoning chain with explicit triple IDs,\\
token budget allocation (two-hop: 40\%, literal: 15\%), faithfulness score: 0.96
}}
\caption{Qualitative reasoning trace demonstrating multi-granular retrieval, token budget allocation, and explicit citation of evidence. The two-hop path and literal are retrieved in parallel based on composite classification, then integrated via structured prompting with constrained decoding.}
\label{fig:qualitative}
\end{figure}

\subsection{Retrieval Efficiency Comparison}
Table \ref{tab:retrieval_speed} contrasts retrieval characteristics. For multi-stage methods we isolate pure retrieval + filtering (excluding generation). Our approach shows markedly lower structural encoding overhead while keeping competitive end-to-end latency due to early granularity pruning.

\begin{table}[t]
\centering
\caption{Retrieval efficiency comparison. Retrieval latency = candidate generation + structural expansion + filtering. End-to-end latency from question input to final answer (reported earlier for baselines). Memory footprint excludes base LLM weights.}
\label{tab:retrieval_speed}
\begin{adjustbox}{max width=\columnwidth}
\begin{tabular}{lcccc}
\toprule
Method & Paradigm & Retrieval Lat. (ms) & End2End (ms) & Mem (GB) \\
\midrule
Text RAG + LLaMA-3 & BM25 + dense passages & 180 & 523 & 2.1 \\
GNN-RAG & GNN-enhanced graph enc. & 420 & 1,050 & 6.4 \\
Think-on-Graph & Iterative expand (multi-step) & 650 & 1,124 & 3.8 \\
KG-Agent & Plan-act loops & 780 & 1,456 & 4.2 \\
Graph-Constrained Reasoning & Lexical + constraint decode & 260 & 1,010 & 2.7 \\
Ours (Zero-shot) & Multi-granular BM25 + hybrid & 290 & 687 & 2.5 \\
Ours (Fine-tuned) & Same retriever + tuned decode & 290 & 592 & 2.5 \\
\bottomrule
\end{tabular}
\end{adjustbox}
\end{table}

\section{Analysis and Discussion}

Our experimental results demonstrate several key insights: (i) Multi-granular decomposition provides targeted retrieval that significantly improves performance across all question types; (ii) Composite label classification enables precise identification of questions requiring multiple knowledge types; (iii) Advanced filtering substantially improves faithfulness while maintaining answer quality; (iv) The framework achieves superior performance while maintaining competitive efficiency compared to existing approaches.

The multi-label classification approach with composite labels proves crucial for system performance, with questions requiring multiple knowledge types showing the largest improvements. This validates our hypothesis that real-world questions often require composite reasoning patterns that cannot be captured by single-label approaches.

\subsection{Comparative Analysis with Referenced Methods}
\textbf{Granularity Selection:} Prior methods either retrieve large undifferentiated neighborhoods (GNN-RAG \cite{mavromatis2024gnn}, Graph-Constrained Reasoning (GCR) \cite{luo2024graph}) or rely on iterative exploration (KG-Agent \cite{jiang2024kg}, Think-on-Graph \cite{sun2023think}). Our multi-label classifier (situated within the taxonomy summarized by \cite{jin2024large}) narrows retrieval \emph{before} expansion, reducing irrelevant edges and aligning with calls for explicit subgraph construction for faithful reasoning \cite{luo2023reasoning}.
\textbf{Retrieval Mechanism:} We employ a lexical-first BM25 layer for triple / literal / description shards with light embedding re-ranking—eliminating GNN message passing latency required in GNN-RAG \cite{mavromatis2024gnn} while preserving semantic sensitivity via secondary cosine scoring.
\textbf{Composite Reasoning Types:} Composite labels (e.g., one-hop+literal) unify evidence attachment in a single pass; iterative / agent frameworks (Think-on-Graph \cite{sun2023think}, KG-Agent \cite{jiang2024kg}) typically trigger additional expansion rounds, and GCR \cite{luo2024graph} focuses on structural validity during decoding rather than pre-retrieval granularity prediction.
\textbf{Filtering Strategy:} Instead of learned graph encoders (e.g., GNN-RAG \cite{mavromatis2024gnn}), we use interpretable structural priors (degree penalty, relation prior, type compatibility) plus diversity regularization, improving faithfulness and traceability consistent with faithful reasoning principles in \cite{luo2023reasoning}.
\textbf{Decoding Constraints:} We combine entity whitelisting and soft relation guidance with citation alignment; GCR \cite{luo2024graph} enforces entity validity but lacks fine-grained granularity-aware evidence budgeting; agent / iterative planners \cite{jiang2024kg,sun2023think} add latency via planning loops.
\textbf{Efficiency Trade-off:} Table \ref{tab:retrieval_speed} shows our retrieval latency sits far below iterative / GNN-heavy methods \cite{mavromatis2024gnn,jiang2024kg,sun2023think} while delivering higher EM / F1 (Table \ref{tab:qa_perf}), validating that early granularity classification plus lexical-first retrieval is an advantageous latency-quality operating point.
\textbf{Maintainability:} BM25 indices and shard inventories can be incrementally updated without re-training structural encoders (contrasting with GNN-based retrievers \cite{mavromatis2024gnn}), simplifying KG refresh cycles.

\section{Limitations and Future Directions}

Despite strong empirical results, our framework has several limitations that point toward future research:

\textbf{KG Coverage and Freshness:} Performance is bounded by KG completeness. Questions about recent events (post-2016, Freebase knowledge cutoff) or long-tail entities with sparse neighborhoods fail. Integration with dynamic KG update pipelines \cite{lehmann2015dbpedia} or hybrid text-KG retrieval \cite{sun2018graftnet} could address this. We note that while BM25-first improves latency, purely lexical anchoring can miss paraphrastic or low-overlap questions; hybrid dense-first variants are future work.

\textbf{Question Classification Errors:} Misclassification (10.7\% error rate on complex paraphrases) cascades downstream—predicting \texttt{one-hop} for a two-hop question retrieves insufficient evidence. Incorporating uncertainty quantification (e.g., conformal prediction \cite{angelopoulos2021gentle}) to trigger fallback retrieval when classifier confidence is low could improve robustness.

\textbf{Two-Hop Limitation:} Current system handles up to two-hop reasoning. Questions requiring three or more hops (4.3\% of ComplexWebQuestions) fail. Extending to $k$-hop requires principled pruning to avoid $O(d^k)$ explosion—promising directions include learned heuristics for intermediate entity selection or hybrid symbolic-neural planning \cite{yao2023react,schick2024toolformer}.

\textbf{Prompt Length Constraints:} The 4096-token LLaMA-3 limit forces truncation for high-degree entities with many neighbors. While our token budget allocation mitigates this, questions requiring broad context (e.g., comparing multiple films of a director) suffer. Future work could explore hierarchical summarization \cite{sarthi2024raptor} or iterative refinement with retrieval-after-initial-answer.

\textbf{Computational Cost:} Despite optimizations, the system requires substantial resources (2 GPU + 32 CPU cores for production throughput). Research into distillation of retrieval models, quantized LLMs, or learned index selection policies could reduce deployment costs.

\textbf{Multilingual and Multi-Modal Extension:} Current implementation is English-only and text-based. Extending to multilingual KGs \cite{lehmann2015dbpedia} or incorporating visual/audio entities (e.g., for film/music questions) requires cross-lingual alignment and multi-modal encoders—significant but feasible extensions of our modular architecture.

\section{Ethical Considerations}

Our focus on factual grounding through knowledge graphs aims to reduce hallucination and improve reliability. However, there is risk of propagating biases present in the underlying knowledge graph. We mitigate this through provenance citation and differential auditing mechanisms. All experiments use publicly available benchmark datasets without processing personally identifiable information.

\section{Conclusion}

We presented a comprehensive framework for knowledge-intensive question answering that addresses fundamental challenges in integrating structured knowledge graphs with large language models. Our approach is grounded in three key insights: (1) \emph{Questions exhibit natural decomposition into distinct knowledge access patterns}—factoid, multi-hop, literal, and definitional—motivating KG partitioning by reasoning depth rather than entity-centric neighborhoods; (2) \emph{Composite reasoning patterns are prevalent} (32-41\% of questions), requiring simultaneous heterogeneous evidence retrieval that single-label or iterative approaches cannot efficiently provide; (3) \emph{Interpretable structural priors} (degree penalties, relation compatibility, type constraints) rival learned graph encoders in precision while offering superior latency, maintainability, and debuggability.

The resulting system combines multi-granular KG decomposition, automated multi-label training data generation via execution trace analysis, hybrid BM25-semantic retrieval with constrained beam expansion, and citation-enforced decoding. Extensive experiments on WebQSP, ComplexWebQuestions, and MetaQA demonstrate state-of-the-art performance: 15.3\% average EM improvement over best baselines, 91.3\% faithfulness score with explicit citation, and sub-800ms end-to-end latency—substantially faster than iterative reasoning methods (1100-1600ms) while maintaining higher accuracy.

Beyond performance gains, our framework offers practical advantages for production deployment: (1) \emph{Modularity}—granularity-specific indices can be updated independently as KGs evolve without retraining neural components, (2) \emph{Transparency}—retrieved evidence directly corresponds to predicted question types, enabling targeted debugging and failure analysis, (3) \emph{Efficiency}—early granularity classification prunes search space before expensive expansion, achieving favorable latency-quality tradeoffs.

Future research directions include: (1) \emph{Adaptive k-hop reasoning} with learned intermediate entity selection to handle questions requiring three or more hops without combinatorial explosion, (2) \emph{Temporal KG integration} for time-sensitive queries through snapshot indexing and temporal consistency checking, (3) \emph{Multi-modal evidence fusion} combining structured KG facts with images, videos, and audio for richer question domains, (4) \emph{Continual learning mechanisms} that incrementally update classifier and retrieval indices as new question patterns emerge, (5) \emph{Cross-lingual extension} leveraging multilingual KGs and language models for non-English question answering.

By explicitly aligning KG structure with question reasoning patterns, our framework demonstrates that thoughtful decomposition and targeted retrieval can unlock the synergies between symbolic knowledge and neural generation—a crucial step toward reliable, explainable, and scalable knowledge-intensive AI systems.

% ----------------------------
% Bibliography
% ----------------------------
\bibliographystyle{IEEEtran}
\bibliography{ref}
\end{document}